% Independently audited by root, PASS, 2026-11-19 16:27 UTC.
\begin{saresult}{A5.014}{Proved under stated hypotheses}{The maximal triviality locus for compact graph fibres}
Let $M$ be an o-minimal expansion of a real closed field, and let
$f:X\to Y$ be a continuous definably proper definable map whose
fibres have dimension at most one.  Then its maximal open
definable-manifold-without-boundary local-triviality locus is
definable, with no restriction on $\dim Y$.
For an arbitrary continuous definably proper map, the part of its
maximal locus over which the fibres have dimension at most one
is also definable.  Combined with A5.013, the remaining question
requires fibre dimension at least two and target dimension at
least three.
\end{saresult}
\begin{proof}
We give a definable construction over the given field; ordinary
compactness of metric spaces and integration of functions are
not used.  Replace $Y$ by its definable manifold locus, which is
open and definable by A5.005.  Embed $X$ in $M^n$, and write
$F_y=f^{-1}(y)$.  Each fibre is a compact finite graph, allowing
isolated points, empty graphs, and disconnected graphs.

\medskip\noindent\textbf{1. Definable arc length and a uniform replacement estimate.}
For a definable embedded arc $\gamma:[a,b]\to M^n$, define its
length by
\[
 \ell(\gamma)=\sum_{j=1}^n\operatorname{Var}_{[a,b]}(\gamma_j).
 \tag{A5.14.1}
\]
Every coordinate has finitely many monotonicity intervals.  If
their number is at most $N$, its variation is the maximum of
\[
 \sum_{i=1}^{N+1}|\gamma_j(t_i)-\gamma_j(t_{i-1})|,
 \qquad a=t_0\leq t_1\leq\cdots\leq t_{N+1}=b.
\]
Repeated partition points are permitted.  This is a maximum on
a compact definable simplex, and is first-order definable.
Length is additive under concatenation, independent of an
increasing reparameterization, and strictly positive on every
nonconstant subarc.  Cumulative length is a continuous strictly
increasing definable function, and hence parameterizes an arc
by an interval $[0,\ell]$.  Ambient distance is bounded by this
length.  A definable embedded circle similarly has a definable
total length $L>0$ and shorter-arc metric $d$; it is the metric
circle of circumference $L$ in these coordinates.

There is a single finite bound $N$ adequate for all the fibres
in the present definable family.  Indeed, Hardt triviality
\cite[Theorem 5.22, pp.62--63]{A5-Coste}, followed by
triangulation of its finitely many compact representative
fibres, gives finitely many definable formula families
$\gamma_e(y,t)$ parameterizing their edges.  Monotonicity and
cell decomposition, with $t$ the last coordinate, give a uniform
bound for every ambient coordinate of every such edge
\cite[Theorems 2.1, 2.9--2.10, pp.15, 19]{A5-Coste}.  Coalescing
bivalent subdivisions multiplies this bound by at most the fixed
number of edges.  Cutting a circle at an arbitrary point adds
at most one break.  Enlarge $N$ once to cover all these cases.

An \emph{intrinsic vertex} of a graph has local valence different
from two; isolated points have valence zero.  The components
left after removing intrinsic vertices are its intrinsic edges.
A component without vertices is a circle.  A homeomorphism
permutes this finite intrinsic graph, is monotone on each
oriented edge, and has one of the two orientations on a circle.
Replace an edge homeomorphism by the map linear in normalized
length coordinates.  On a circle prescribe the image of one
base point and make the same replacement around the circle.
Call these maps \emph{canonical}.

Suppose $h:F_y\to F_z$ is a definable homeomorphism satisfying
$\|h(x)-x\|<\epsilon$ for all $x$.  On one source edge, transport
the target coordinate functions along $h$.  Each coordinate of
both curves has at most $N$ monotonicity intervals.  Refining
their two partitions gives at most $2N+2$ intervals.  On each
interval the absolute difference between the two coordinate
variations is at most $2\epsilon$, by comparing the two
endpoints.  Thus the source prefix length $u$ and target prefix
length $v$ of $h(x)$ satisfy
\[
       |u-v|\leq C\epsilon,
       \qquad |\ell_y-\ell_z|\leq C\epsilon,
       \qquad C=4n(N+1).
 \tag{A5.14.2}
\]
The canonical replacement sends $u$ to
$v'=u\ell_z/\ell_y$.  Since $0\leq u\leq\ell_y$,
$|v'-u|\leq C\epsilon$, so $|v'-v|\leq2C\epsilon$.
Ambient distance is bounded by length.  Consequently its
displacement is less than
\[
                         (1+2C)\epsilon.                 \tag{A5.14.3}
\]
The argument applies to a cut circle with base point image
$h(p)$.  For a loop based at an intrinsic vertex, use the
interval parametrization whose two endpoints are identified
at that vertex; the same variation and prefix-length estimates
apply to this cut parametrization.  At isolated vertices
the original displacement estimate suffices.

All canonical maps occur in finitely many first-order formula
formats.  The finite data specify the Hardt pieces of $y,z$,
an isomorphism of their intrinsic graphs, and the orientations
of the edges.  Only the images of the chosen circle base points
are additional geometric parameters.  Formula (A5.14.1) and
the unique inverse of cumulative length define the map on each
edge; the definitions agree at the vertices.  For circles use
length modulo their circumference.  Thus the graph, its
homeomorphism property, and its displacement bound are all
expressible by finitely many formulas.  This conclusion uses
the one-dimensional monotonicity estimate, not an assumed
uniform bound for arbitrary graph homeomorphisms.

\medskip\noindent\textbf{2. Complete regularity is definable.}
Call $y$ completely regular if for every $\epsilon>0$ there is
$\delta>0$ such that every $z\in Y$ with $\|z-y\|<\delta$
has a definable homeomorphism $h:F_y\to F_z$ of displacement
less than $\epsilon$.  Step~1 shows that the same condition is
obtained by requiring $h$ to be canonical: use
$\epsilon/(1+2C)$ before applying (A5.14.3).  Its precise
first-order form is therefore
\[
 \begin{gathered}
 \forall\epsilon>0\ \exists\delta>0\ \forall z\in Y,\quad
 \|z-y\|<\delta\ \Longrightarrow\\
 \bigvee_i\exists b_i\,
 \left(C_i(y,z,b_i)\ \wedge\
 \forall(x,x')\in\Gamma_{i,y,z,b_i},\ \|x'-x\|<\epsilon\right),
 \end{gathered}                                           \tag{A5.14.4}
\]
where $C_i$ specifies a canonical homeomorphism and
$\Gamma_i$ its graph.  Let $R_f\subset Y$ be the definable set
given by (A5.14.4).

If $f$ is locally trivial at $y$, make its product chart the
identity on $F_y$.  Uniform continuity of this chart on the
product of $F_y$ with a compact small target neighbourhood
gives (A5.14.4).  The same applies at every point of that
neighbourhood.  Hence the maximal locus is contained in
$\operatorname{Int}_Y R_f$.  It remains to prove the converse.

\medskip\noindent\textbf{3. The metric circles vary continuously on $R_f$.}
Fix a point $y_0\in\operatorname{Int}R_f$ and write $F_0=F_{y_0}$.
The empty-fibre case is immediate from (A5.14.4), which then
forces all nearby fibres to be empty.  Otherwise shrink to an
open neighbourhood $V$ contained in $\operatorname{Int}R_f$
such that every fibre is canonically homeomorphic to $F_0$
by a map of a fixed sufficiently small displacement.  Distinct
components of $F_0$ have positive separation.  Thus its circle
components can be labelled throughout $V$ by disjoint fixed
ambient neighbourhoods.  Their unions in $f^{-1}(V)$ are both
open and closed relative to $f^{-1}(V)$, since the displacement
bound leaves a positive gap between those neighbourhoods and
the other components.

For one such circle family $C_y$, denote its length and metric
by $L_y,d_y$.  They are definable by Step~1.  They are also
continuous, with $d_y$ jointly continuous on the fibre-pair
family.  In fact, at any $y\in V$, complete regularity provides
canonical comparisons $k:C_y\to C_z$ with arbitrarily small
displacement as $z\to y$.  Equation (A5.14.2) shows
$L_z\to L_y$.  Since the comparison is linear in normalized
length,
\[
 d_z(k(x),k(x'))=\frac{L_z}{L_y}d_y(x,x').                 \tag{A5.14.5}
\]
For arbitrary $x_z,x'_z$ approaching $x,x'$, their inverse
images under $k$ approach $x,x'$ in the fixed compact circle
$C_y$, because $k^{-1}$ has the same displacement bound as
$k$.  The metric $d_y$ has the ordinary definable topology
on this fixed circle and is uniformly continuous.  Equation
(A5.14.5) proves the assertion by epsilon-delta estimates.
This reasoning is valid in the given ordered field; no
sequential compactness is being substituted for definable
compactness.

\medskip\noindent\textbf{4. Continuous definable anchors for the circle components.}
Choose a regular $C^1$ point $p\in C_0$; such points exist by
definable $C^1$ cell decomposition
\cite[Theorem 6.6, p.71]{A5-Coste}.  There is a constant
$A>0$ such that
\[
                     d_0(p,z)\leq A\|z-p\|\qquad(z\in C_0).
 \tag{A5.14.6}
\]
Locally at $p$, choose one coordinate with nonzero tangent
derivative and use it as the arc parameter.  The other
coordinate derivatives are bounded on a smaller compact arc,
so its $L^1$ length is bounded by a constant times that
coordinate difference.  Away from the arc, compactness gives
a positive lower bound for $\|z-p\|$, whereas $d_0\leq L_0/2$.
These two estimates prove (A5.14.6).

For a constant $M>A^2$ to be enlarged below, define
\[
 E_y(x)=\max_{z\in C_y}
             \left(d_y(x,z)^2-M\|z-p\|^2\right).
 \tag{A5.14.7}
\]
The maximum exists and is definable.  By (A5.14.6),
$E_0(p)=0$.  The term $z=p$ gives
$E_0(x)\geq d_0(x,p)^2$, so $p$ is its unique global
minimizer.  Under a small canonical comparison
$k:C_0\to C_y$, (A5.14.5), $L_y\to L_0$, and
$\sup\|k(z)-z\|\to0$ show that
$E_y(k(x))\to E_0(x)$ uniformly in $x$.  This follows directly
by bounding the difference inside the maximum in (A5.14.7);
all lengths and ambient distances are uniformly bounded on
the compact reference circle and its small neighbourhood.

Choose concentric short closed intrinsic arcs $J_0\subset
\operatorname{Int}I_0$ about $p$, with the radius of $I_0$
less than $L_0/8$, and $J_0$ strictly smaller.  Points of
$C_0\setminus\operatorname{Int}I_0$ have ambient distance
at least $m>0$ from $p$.  Enlarge $M$ so that, for all nearby
circles, a point at distance at least $m/2$ from $p$ has
\[
            d_y(x,z)^2-M\|z-p\|^2<-1
\]
for every $x\in C_y$; use a fixed upper bound for $L_y/2$.
The point $k(p)$, on the other hand, contributes a value
greater than $-1/2$ when $y$ is sufficiently close to $y_0$.
Thus every maximizer defining (A5.14.7) lies in $I_y=k(I_0)$.
The strict minimum and positive gap of $E_0$ outside
$\operatorname{Int}J_0$, together with the preceding uniform
estimate, show that every global minimizer of $E_y$ lies in
$\operatorname{Int}J_y$, where $J_y=k(J_0)$.

For $x\in J_y$ and $z\in I_y$ the shorter arc has no
wrap-around.  In an intrinsic length coordinate $u$ on
$I_y$, equation (A5.14.7) becomes
\[
 E_y(u)=u^2+\max_{z\in I_y}
       \left(-2u\,u(z)+u(z)^2-M\|z-p\|^2\right).
 \tag{A5.14.8}
\]
The second term is convex as a maximum of affine functions.
More explicitly, for $0<t<1$ its values satisfy
\[
 E_y((1-t)u+tv)\leq(1-t)E_y(u)+tE_y(v)-t(1-t)(u-v)^2.
\]
Consequently $E_y$ is strictly convex on $J_y$.  Since all
global minimizers lie there, there is exactly one; denote it
by $s(y)$.  The equation of a global minimizer is a first-order
condition, so $s$ is definable.

It is continuous throughout a smaller $V$.  At any $y\in V$,
small canonical comparisons and (A5.14.5) make the energies
uniformly close on the fixed reference circle $C_y$.  That
circle has a unique minimizer of its continuous energy.
Outside any chosen neighbourhood of that minimizer, compactness
gives a positive gap above the minimum.  Uniform closeness
therefore puts the minimizer for every sufficiently close
$z$ in the corresponding neighbourhood.  Small displacement
of the comparison makes this an ambient neighbourhood as well.
This proves continuity by epsilon-delta estimates.  Repeating
the construction for the finitely many circle components
supplies continuous definable anchors $s_j(y)$, with
$s_j(y_0)=p_j$.

\medskip\noindent\textbf{5. Unique small canonical maps produce a bundle.}
Mark $F_0$ at the chosen point $p_j$ on each circle component.
Its canonical automorphisms fixing these marked points form
a finite group $G$: the intrinsic finite graph has only
finitely many permutations and orientations; on a marked
circle there are just the identity and reflection.  Choose
$\eta>0$ so that
\[
 2\eta<\min_{a\in G,\ a\ne1}\max_{x\in F_0}\|a(x)-x\|,
 \tag{A5.14.9}
\]
and choose any $\eta>0$ if $G$ is trivial.  The positive
minimum exists because the set is finite and every nonidentity
map moves some point.

After shrinking $V$, there exists for every $y\in V$ a
canonical homeomorphism $h_y:F_0\to F_y$ with
$h_y(p_j)=s_j(y)$ and displacement less than $\eta$.
To see this, begin with an arbitrarily small canonical map
given by complete regularity at $y_0$.  On each target circle,
rotate its length coordinate to send the image of $p_j$ to
$s_j(y)$.  The required intrinsic rotation tends to zero:
transport back to the fixed circle by (A5.14.5), and use
continuity of its metric and of $s_j$.  Ambient motion is
at most the intrinsic rotation.  The adjusted canonical map
therefore has displacement less than $\eta$.

It is unique.  If $h,h'$ were two such maps, then
$a=h^{-1}h'\in G$, and for every $x\in F_0$,
\[
 \|a(x)-x\|\leq\|a(x)-h(a(x))\|+\|h'(x)-x\|<2\eta.
\]
Equation (A5.14.9) forces $a=1$.  The finite formula formats
from Step~1, together with the anchor conditions and the
strict displacement bound, consequently define the graph of
a single family $h_y$.

This family is continuous.  Fix $y\in V$.  Complete regularity
at $y$ gives canonical comparisons $k:F_y\to F_z$ of
arbitrarily small displacement for $z$ near $y$.  For small
enough displacement the labelled circle components match.
Rotate them as above to make $k(s_j(y))=s_j(z)$; by continuity
of the anchors and (A5.14.5), this changes the displacement
by a quantity tending to zero.  Call the result $k'$.  The
composite $k'h_y$ is canonical, respects all anchors, and has
displacement below $\eta$ for $z$ sufficiently close to $y$,
since $h_y$ has a strict margin below $\eta$.  Uniqueness gives
\[
                             h_z=k'h_y.
\]
Hence $h_z\to h_y$ uniformly in the ambient metric.  This
implies joint continuity of
\[
 H:F_0\times V\longrightarrow f^{-1}(V),\qquad
 H(x,y)=h_y(x),\qquad fH=\operatorname{pr}_V.
 \tag{A5.14.10}
\]
The map is a definable continuous bijection.  It is definably
proper: the preimage of a compact set $A\subset f^{-1}(V)$
is a closed subset of the compact set $F_0\times f(A)$.
The definable closed-image argument from A5.012 makes its
inverse continuous.  Thus (A5.14.10) is a local definable
bundle chart.

We have proved that the maximal locus is exactly
$\operatorname{Int}_Y R_f$, which is definable.  Finally,
for an arbitrary proper map let
$Z=\{y:\dim F_y\leq1\}$.  A locally trivial value with such
a fibre lies in $\operatorname{Int}_Y Z$.  Restrict the map
to that definable open set and apply the result just proved.
This gives precisely the claimed low-dimensional-fibre part
of the maximal locus and proves the final reduction.
\end{proof}
